% Generated draft: requires mathematical and source review.
\documentclass[11pt]{article}
\usepackage[margin=25mm]{geometry}
\usepackage{fontspec}
\setmainfont{Latin Modern Roman}
\usepackage{amsmath,amssymb,mathtools,unicode-math}
\setmathfont{Latin Modern Math}
\usepackage{enumitem,needspace,xurl,hyperref}
\setlength{\parindent}{0pt}
\setlength{\parskip}{5pt}
\setlength{\emergencystretch}{3em}
\newcommand{\N}{\mathbb{N}}
\newcommand{\Z}{\mathbb{Z}}
\newcommand{\Q}{\mathbb{Q}}
\newcommand{\R}{\mathbb{R}}
\newcommand{\C}{\mathbb{C}}
\newcommand{\F}{\mathbb{F}}
\newcommand{\A}{\mathbb{A}}
\newcommand{\PP}{\mathbb{P}}
\newcommand{\E}{\mathbb{E}}
\newcommand{\eps}{\varepsilon}
\newcommand{\dd}{\,\mathrm{d}}
\begin{document}
\section*{UnsolvedMath ID: 30005998}
\textbf{Source problem code:} catalogue entry
\subsection{Definitions and mathematical statement}
The supplied record asks for a classification of entire stable Allen--Cahn configurations in \(\R^3\). It does not specify the potential, the admissible class of configurations, or the precise macroscopic conclusions sought. The following gives a conditional variational framework without supplying those missing hypotheses.

In a scalar formulation, let \(W\in C^2(\R)\) denote the intended potential and let \(\mathcal U\subseteq C^2(\R^3;\R)\) denote the intended admissible class. Both must be specified to obtain a definite classification problem. For a bounded open set \(\Omega\subset\R^3\), define the local energy
\[
 E_W(u;\Omega)=\int_\Omega\left(\frac12|\nabla u|^2+W(u)\right)\,\dd x.
\]
An entire critical point satisfies
\[
 -\Delta u+W'(u)=0\qquad\text{in }\R^3.
\]
Under the usual compactly supported variation convention, it is stable if
\[
 \int_{\R^3}\left(|\nabla\varphi|^2+W''(u)\varphi^2\right)\,\dd x\geq0
 \qquad\text{for every }\varphi\in C_c^\infty(\R^3).
\]
Thus the conditional classification target is
\[
 \mathcal S(W,\mathcal U)=\left\{u\in\mathcal U:
 -\Delta u+W'(u)=0\text{ in }\R^3,\quad
 \int_{\R^3}\left(|\nabla\varphi|^2+W''(u)\varphi^2\right)\,\dd x\geq0
 \text{ for every }\varphi\in C_c^\infty(\R^3)
 \right\}.
\]
For comparison, an absolute minimizer under compact perturbations satisfies
\[
 E_W(u;\Omega)\leq E_W(u+\psi;\Omega)
\]
for every bounded open \(\Omega\subset\R^3\) and every \(\psi\in C_c^\infty(\Omega)\), subject to any additional admissibility restrictions in the intended formulation. These local definitions do not require finite total energy on \(\R^3\).

\textbf{Requested problem, with unresolved qualifications.} Specify the intended \(W\) and \(\mathcal U\), and classify every member of \(\mathcal S(W,\mathcal U)\). In particular, determine whether the macroscopic regularity and minimal-surface behavior referred to in the source for absolute minimizers extend to every stable configuration. The supplied excerpt does not identify those results or define their scaling, topology of convergence, limiting objects, or regularity conclusions. Consequently this last request cannot yet be expressed as a precise quantified assertion. The framework above does not assert that the problem is posed for every \(C^2\) potential or every entire scalar solution.
\subsection{Short English statement}
The problem asks which entire Allen--Cahn configurations in \(\R^3\) have nonnegative second variation, and whether they share the large-scale regularity and minimal-surface behavior attributed to absolute minimizers. The potential, admissibility assumptions, and exact large-scale conclusions remain unspecified. The supplied record labels the problem open; this draft does not independently verify that assessment.
\subsection{Sources}
\begin{itemize}
\item Supplied local UnsolvedMath registry record, permanent ID \(30005998\), title \emph{Classification of Stable Allen--Cahn Configurations}, source problem code \texttt{catalogue entry}; collection: \emph{Oberwolfach Reports Open Problems}. This draft uses the supplied record and does not claim an independent inspection of repository files.
\item Supplied report citation: \emph{Calculus of Variations (2025)}, Oberwolfach Reports, DOI \texttt{10.4171/owr/2024/37}, pp.\ 2113--2178. \url{https://doi.org/10.4171/owr/2024/37}. The citation is retained as supplied, including its year designation.
\item Attribution in the supplied excerpt: Joaquim Serra, \emph{Stable phase transitions: open questions and new results}, joint work with Hardy Chan, Alessio Figalli, and Xavier Fern\'andez-Real. The excerpt asks about extending preceding results to all stable configurations, informally described there as those observable in nature; the preceding results themselves are not supplied. No separate proposer is identified.
\item Supplied literature reference: Chan--Fernandez-Real--Figalli--Serra, \url{https://arxiv.org/abs/2503.21245}. No paper title or precise theorem statement is supplied.
\item Supplied literature reference: \emph{Oberwolfach Report 37/2024}, \url{https://ems.press/content/serial-article-files/50045}.
\item Literature assessment dated exactly \texttt{2026-08-22}: the supplied curation labels the problem open and states that results for minimizers and selected stable configurations do not classify all stable configurations. This is supplied dated curation, not an independently verified literature assessment; the background explicitly requests independent verification of time-sensitive status and cited claims.
\item The registry labels the statement \texttt{corrected\_verified} and says it was reconstructed and checked against the official Oberwolfach report. This is a supplied verification claim, not verification performed for this draft.
\end{itemize}
\end{document}
